\subsubsection{Pollard's Rho + factorizacion}
\begin{lstlisting}[style=mystyle, language=C++]
// TC: expected O(n^{1/4}) per factor
// usa: descomponer numeros hasta 2^64 en factores primos
#include <bits/stdc++.h>
using namespace std;

using u64 = unsigned long long;
using u128 = __uint128_t;

// assumes millerRabin and modpow_mr are defined above

static u64 modmul(u64 a, u64 b, u64 mod){
	return (u128)a * b % mod;
}

static u64 pollard(u64 n){
	if(n % 2 == 0) return 2;
	static std::mt19937_64 rng((u64)chrono::steady_clock::now().time_since_epoch().count());
	while(true){
		u64 c = uniform_int_distribution<u64>(1, n - 1)(rng);
		u64 x = uniform_int_distribution<u64>(0, n - 1)(rng);
		u64 y = x;
		u64 d = 1;
		auto f = [&](u64 v){ return (modmul(v, v, n) + c) % n; };
		while(d == 1){
			x = f(x);
			y = f(f(y));
			d = std::gcd<u64>(x > y ? x - y : y - x, n);
		}
		if(d != n) return d;
	}
}

// returns sorted list of prime factors (with multiplicity)
void factorRec(u64 n, vector<u64>& out){
	if(n == 1) return;
	if(millerRabin(n)){ out.push_back(n); return; }
	u64 d = pollard(n);
	factorRec(d, out);
	factorRec(n / d, out);
}

vector<u64> factorize(u64 n){
	vector<u64> out;
	factorRec(n, out);
	sort(out.begin(), out.end());
	return out;
}

int main(){
	// factor 123456789012345 = 3^2 * 5 * 7 * 13 * 29 * 41 * 271 * 9091
	auto f = factorize(123456789012345ULL);
	for(u64 p : f) cout << p << " ";
	cout << "\n";
	return 0;
}
\end{lstlisting}
